\documentclass{article}
\usepackage{amsmath}
\usepackage{cancel}
\usepackage{graphicx} % for creating figures

\title{States in Thermal Physics}

\begin{document}
\maketitle
The root of statistical mechanics is the concept of the quantum mechanical
stationary states of a system. Recall that stationary states are defined with
respect to the energy eigenstates of the Hamiltonian, i.e., states satisfying
\begin{equation}
  \hat{H}|\varphi_n\rangle = E_n|\varphi_n\rangle
\end{equation}

Such a state is considered \textit{stationary} because it's time dependence is
given by
\begin{equation}
  |\varphi_n(t)\rangle = e^{-iE_n t/\hbar}|\varphi_n\rangle
\end{equation}
In other words, the time dependence is a negligible \textit{overall phase}; for
all intents and purposes, there is no detectable time dependence.

As we can see in the above equation, each of these stationary states has a
definite energy. If many states share the same energy, we call that
energy value \textit{degenerate}. For example a hydrogen atom in the $n=2$
energy state has energy $-13.6/4$ eV. Such a state has an allowed total angular
momentum quantum number of $\ell = 1$ but there is \textit{more than one
  possibility} for the z-projection of the angular momentum, namely $m_\ell =
-1, \; 0, \; 1$. This means that the states
\begin{equation}
  |2,1,1\rangle, \qquad |2,1,0\rangle, \qquad |2,1,-1\rangle
\end{equation}
are degenerate with a \textit{threefold degeneracy}. Another term for the
degeneracy of a particular energy value is the \textit{multiplicity}. For a
system of many particles each with its own energy spectrum, we must be careful
to keep track of the possible degeneracies. For example, a system of
three non-interacting harmonic oscillators in one dimension has an energy
spectrum given by
\begin{equation}
  E_{n_1,n_2,n_3} = \frac{\hbar^2\pi^2}{2mL^2}(n_1^2+n_2^2+n_3^2)
\end{equation}
where I have chosen to shift the $|0,0,0\rangle$ state energy to a value of $0$.
The first excited state of the system can be made in any of the following ways
\begin{equation}
  |1,0,0\rangle, \quad |0,1,0\rangle, \quad |0,0,1\rangle
\end{equation}
and therefore, we identify the energy $\frac{\hbar^2\pi^2}{2mL^2}$ as having a
multiplicity of $3$. \textit{Note:} this is the same energy spectrum as for the
3-dimensional isotropic harmonic oscillator.



\end{document}






























